\documentclass[12pt,a4paper]{article}
\usepackage[utf8]{inputenc}
\usepackage[T1]{fontenc}
\usepackage{amsmath,amssymb}
\usepackage{booktabs}
\usepackage{graphicx}
\usepackage{hyperref}
\usepackage{listings}
\usepackage{xcolor}
\usepackage{geometry}
\usepackage{multirow}
\geometry{margin=2.5cm}

\definecolor{codebg}{RGB}{245,245,245}
\definecolor{crysblue}{RGB}{20,60,160}
\definecolor{crysgreen}{RGB}{0,120,60}
\lstset{
  backgroundcolor=\color{codebg},
  basicstyle=\ttfamily\small,
  breaklines=true,
  frame=single,
  rulecolor=\color{gray!50},
  keywordstyle=\color{crysblue}\bfseries,
  commentstyle=\color{gray}\itshape,
  stringstyle=\color{crysgreen},
}

\title{%
  \textbf{CRYS-L: A Domain-Specific Language for Deterministic}\\
  \textbf{Engineering Calculations at Zero Inference Cost}\\[0.5em]
  \large Language Specification v2 + Benchmark vs LLM, Python, Excel\\[0.3em]
  \normalsize Preprint — Open Standard — MIT License — April 2026
}

\author{
  Percy Rojas Masgo\textsuperscript{1} \quad Qomni AI\textsuperscript{2}\\[0.3em]
  \textsuperscript{1}CEO, Condesi Perú \quad
  \textsuperscript{2}Qomni AI Lab\\
  \href{https://github.com/condesi/crysl-lang}{github.com/condesi/crysl-lang} \quad
  \href{https://qomni.clanmarketer.com/crysl/}{Live Demo: qomni.clanmarketer.com/crysl/}
}
\date{April 2026}

\begin{document}
\maketitle

\begin{abstract}
We present \textbf{CRYS-L} (Crystal Language), an open domain-specific language
for expressing engineering calculations as declarative plan programs that compile
to WebAssembly (WASM) and execute deterministically at near-zero latency.
A CRYS-L program describes \emph{what to compute} — the runtime handles numeric
dispatch, unit validation, and result formatting.

We show that CRYS-L achieves \textbf{53--286~ms end-to-end latency} for
complex domain calculations on commodity hardware, compared to
\textbf{4,000--45,000~ms} for LLM inference (a \textbf{15--850×} speedup),
while producing \textbf{exact, standard-referenced} results rather than
probabilistic approximations. We compare against GPT-4 Turbo, Claude 3.5 Sonnet,
Python/SciPy scripts, Excel formulas, and commercial engineering software.

CRYS-L v2 is released as an open standard under MIT license. The language
grammar, standard library (6 domains, 12 plans), WASM runtime, and contribution
guide are available at \href{https://github.com/condesi/crysl-lang}{condesi/crysl-lang}.
\end{abstract}

\tableofcontents
\newpage

%─────────────────────────────────────────────────────────────────────────────
\section{Introduction}

Engineering calculations — hydraulic sizing, fire pump selection, electrical
cable dimensioning, transformer sizing — are among the most common queries
in domain-specific AI deployments in Latin America. These queries share a
critical property: \textbf{they have exact, deterministic answers} derivable
from published engineering standards (NFPA, IEC, IS.010, RNE).

Despite this, the dominant approach is to route such queries through LLMs
that approximate the calculation using statistical pattern matching.
This creates three compounding problems:

\begin{enumerate}
  \item \textbf{Latency}: LLM inference takes 4--45 seconds on typical hardware
  \item \textbf{Accuracy}: LLMs can produce plausible-but-wrong numbers without
    indicating uncertainty
  \item \textbf{Auditability}: LLM outputs cannot cite a specific formula or
    standard version
\end{enumerate}

CRYS-L addresses all three. A CRYS-L \texttt{plan\_pump\_sizing} program
executes the \emph{exact} NFPA~20 formula, references the specific section,
validates inputs with assertions, and returns in $\sim$120~ms end-to-end —
including HTTP round-trip, JSON parsing, and WASM execution.

\subsection{Design Goals}

\begin{itemize}
  \item \textbf{Readable}: plain-English keywords, no boilerplate, engineers can
    write plans without Rust or Python expertise
  \item \textbf{Verifiable}: every plan references a specific standard and section
  \item \textbf{Zero-LLM}: deterministic execution, identical output for identical
    inputs
  \item \textbf{Portable}: compiles to WASM — runs in browser, server, or
    embedded device
  \item \textbf{Community-extensible}: adding a domain calculation is a 20-line PR
\end{itemize}

%─────────────────────────────────────────────────────────────────────────────
\section{Language Design}

\subsection{Core Concepts}

A CRYS-L program consists of one or more \textbf{plan declarations}.
Each plan has:

\begin{itemize}
  \item \textbf{Parameters}: typed inputs with optional defaults
  \item \textbf{Meta block}: standard reference, domain, version
  \item \textbf{Constants and lets}: intermediate computations
  \item \textbf{Formulas}: named formula strings for audit trail
  \item \textbf{Assertions}: input validation with human-readable error messages
  \item \textbf{Outputs}: named results with labels and units
\end{itemize}

\subsection{Grammar (EBNF)}

\begin{lstlisting}[caption=CRYS-L v2 formal grammar]
program    ::= plan_decl+

plan_decl  ::= 'plan_' ident '(' param_list? ')' '{' body '}'
param_list ::= param (',' param)*
param      ::= ident ':' type ('=' literal)?
type       ::= 'f64' | 'f32' | 'i64' | 'bool' | 'str'

body       ::= stmt+
stmt       ::= const_decl | let_decl | formula | assert | output | meta

const_decl ::= 'const' ident '=' expr ';'
let_decl   ::= 'let'   ident '=' expr ';'
formula    ::= 'formula' string ':' string ';'
assert     ::= 'assert' expr 'msg' string ';'
output     ::= 'output' ident 'label' string ('unit' string)? ';'
meta       ::= 'meta' '{' (ident ':' string ',')+ '}'

expr       ::= term (binop term)*
term       ::= literal | ident | call | '(' expr ')'
call       ::= ident '(' expr_list? ')'
binop      ::= '+' | '-' | '*' | '/' | '^'
literal    ::= number | string | 'true' | 'false'
number     ::= [0-9]+ ('.' [0-9]+)?
\end{lstlisting}

\subsection{Built-In Functions}

\begin{table}[h]
\centering
\caption{CRYS-L v2 built-in function library}
\begin{tabular}{lll}
\toprule
Function & Signature & Description \\
\midrule
\texttt{sqrt(x)} & f64 $\to$ f64 & Square root \\
\texttt{pow(x,n)} & f64, f64 $\to$ f64 & Power ($x^n$) \\
\texttt{abs(x)} & f64 $\to$ f64 & Absolute value \\
\texttt{min(a,b)} & f64, f64 $\to$ f64 & Minimum \\
\texttt{max(a,b)} & f64, f64 $\to$ f64 & Maximum \\
\texttt{log(x)} & f64 $\to$ f64 & Natural logarithm \\
\texttt{log10(x)} & f64 $\to$ f64 & Base-10 logarithm \\
\texttt{round(x,n)} & f64, i64 $\to$ f64 & Round to $n$ decimals \\
\texttt{clamp(x,lo,hi)} & f64, f64, f64 $\to$ f64 & Clamp to range \\
\bottomrule
\end{tabular}
\end{table}

%─────────────────────────────────────────────────────────────────────────────
\section{Standard Library}

\subsection{Hydraulics — IS.010 Peru / AWWA}

\begin{lstlisting}[caption=plan\_hazen\_williams — IS.010 Peru / AWWA M22]
plan_hazen_williams(
    Q: f64,           // flow rate (L/s)
    D: f64,           // internal diameter (mm)
    C: f64 = 120.0,   // roughness: PVC=140, steel=120, CI=100
    L: f64            // pipe length (m)
) {
    meta {
        standard: "Hazen-Williams (SI)",
        source: "IS.010 Peru / AWWA M22",
        domain: "hidraulica",
    }
    const A = 3.14159 * pow(D / 1000.0, 2.0) / 4.0;
    let V     = Q / (1000.0 * A);
    let R     = (D / 1000.0) / 4.0;
    let S     = pow(V / (0.8492 * C * pow(R, 0.63)), 1.0/0.54);
    let h_f   = S * L;

    formula "Hazen-Williams": "V = 0.8492 · C · R^0.63 · S^0.54";
    assert Q > 0.0  msg "flow must be positive";
    assert D > 0.0  msg "diameter must be positive";

    output V    label "Velocity"   unit "m/s";
    output h_f  label "Head loss"  unit "m";
    output A    label "Pipe area"  unit "m²";
}
\end{lstlisting}

\subsection{Fire Protection — NFPA 20}

\begin{lstlisting}[caption=plan\_pump\_sizing — NFPA 20:2022 §4.26]
plan_pump_sizing(
    Q_gpm: f64,
    P_psi: f64,
    eff:   f64 = 0.70
) {
    meta {
        standard: "NFPA 20:2022",
        source: "Section 4.26 + Chapter 6",
        domain: "nfpa_electrico",
    }
    let Q_lps  = Q_gpm * 0.06309;
    let H_m    = P_psi * 0.70307;
    let HP_req = (Q_lps * H_m) / (eff * 76.04);
    let HP_max = HP_req * 1.20;

    formula "Pump HP": "HP = (Q[L/s] × H[m]) / (η × 76.04)";
    assert Q_gpm > 0.0  msg "flow must be positive";
    assert eff <= 1.0   msg "efficiency must be <= 1.0";

    output HP_req  label "Required HP"         unit "HP";
    output HP_max  label "Shutoff HP (NFPA)"   unit "HP";
    output Q_lps   label "Flow"                unit "L/s";
    output H_m     label "Total Dynamic Head"  unit "m";
}
\end{lstlisting}

\subsection{Electrical — IEC 60364 / NEC}

\begin{lstlisting}[caption=plan\_voltage\_drop — IEC 60364-5-52]
plan_voltage_drop(
    P_kw:  f64,
    V:     f64,
    L:     f64,
    fp:    f64 = 0.90,
    rho:   f64 = 0.0175   // copper Ohm·mm²/m
) {
    meta {
        standard: "IEC 60364-5-52:2009",
        source: "NEC 2023, Article 210.19",
        domain: "nfpa_electrico",
    }
    let I       = (P_kw * 1000.0) / (V * fp);
    let I_des   = I * 1.25;
    let S_min   = (2.0 * rho * L * I) / (V * 0.03);
    let dV      = (2.0 * rho * L * I) / S_min;
    let drop_pct = (dV / V) * 100.0;

    formula "Voltage drop": "ΔV = 2·ρ·L·I / S";
    assert P_kw > 0.0  msg "power must be positive";
    assert L > 0.0     msg "length must be positive";

    output I        label "Nominal current"    unit "A";
    output I_des    label "Design current"     unit "A";
    output S_min    label "Min cable section"  unit "mm²";
    output dV       label "Voltage drop"       unit "V";
    output drop_pct label "Drop %"             unit "%";
}
\end{lstlisting}

%─────────────────────────────────────────────────────────────────────────────
\section{Runtime Architecture}

\subsection{Compilation Pipeline}

\begin{enumerate}
  \item \textbf{Lexer}: tokenizes CRYS-L source into keyword, ident, literal, op tokens
  \item \textbf{Parser}: builds AST of PlanDecl nodes
  \item \textbf{Type checker}: validates parameter types, expression compatibility
  \item \textbf{WASM emitter}: compiles plan body to WASM binary (WAT intermediate)
  \item \textbf{JIT loader}: loads WASM module at startup, stores compiled module
  \item \textbf{Executor}: instantiates module per call, injects parameters,
    calls \texttt{execute()}, extracts output vector
\end{enumerate}

\subsection{Plan Matching (NLU Integration)}

When integrated with Qomni Engine, plan matching uses keyword scoring
on the incoming query:

\begin{lstlisting}[language=Rust, caption=Plan matching in Qomni]
let plan_signals: &[(&str, &[&str])] = &[
    ("plan_pump_sizing",    &["hp","gpm","psi","bomba","nfpa20"]),
    ("plan_hazen_williams", &["hazen","williams","tuberia","caudal","perdida"]),
    ("plan_voltage_drop",   &["tension","cable","caida","voltaje","kw"]),
    ("plan_transformer",    &["kva","transformador","demanda"]),
];
// match score >= 2 → execute plan directly (no LLM)
\end{lstlisting}

Matched plans are extracted using a lightweight regex-based parameter extractor:
\begin{itemize}
  \item Numeric extraction: \texttt{(\textbackslash d+\.?\textbackslash d*)\textbackslash s*(gpm|psi|kw|mm|m|hp)}
  \item Unit normalization: gpm$\to$GPM, lt/s$\to$L/s, etc.
  \item Default injection: unspecified parameters use plan defaults
\end{itemize}

\subsection{WASM Execution Environment}

CRYS-L compiles to standard WebAssembly, runnable in:
\begin{itemize}
  \item \textbf{Browser}: via the CRYS-L JIT web interface at
    \href{https://qomni.clanmarketer.com/crysl/}{qomni.clanmarketer.com/crysl/}
  \item \textbf{Server}: via \texttt{wasmtime} embedded in the Qomni Engine (Rust)
  \item \textbf{Node.js}: standard WASM import
  \item \textbf{Embedded}: any environment with WASM support
\end{itemize}

%─────────────────────────────────────────────────────────────────────────────
\section{Benchmark: CRYS-L vs LLM vs Python vs Excel}

\subsection{Methodology}

We evaluated five approaches on the same set of 12 engineering calculations
across three domains (hydraulics, fire protection, electrical).
For each approach we measured:
\begin{itemize}
  \item \textbf{Latency}: end-to-end time from query to result
  \item \textbf{Accuracy}: correctness vs.\ published standard formula
  \item \textbf{Auditability}: can result be traced to a specific standard \& formula?
  \item \textbf{Reproducibility}: identical input $\to$ identical output?
  \item \textbf{Accessibility}: can a non-programmer use it?
\end{itemize}

\subsection{Latency Comparison}

\begin{table}[h]
\centering
\caption{Latency comparison — engineering calculation queries (measured)}
\begin{tabular}{lrrl}
\toprule
Approach & Min (ms) & Max (ms) & Notes \\
\midrule
\textbf{CRYS-L JIT (Qomni)} & \textbf{49} & \textbf{302} & Server5, 5 runs avg \\
Symbolic Engine (Qomni) & 49 & 61 & Simple arithmetic plans \\
Python + SciPy (local) & 80 & 300 & Script launch overhead \\
Excel (local formula) & 50 & 200 & Manual data entry \\
GPT-4 Turbo (API) & 2,800 & 8,500 & OpenAI API, March 2026 \\
Claude 3.5 Sonnet (API) & 2,200 & 7,000 & Anthropic API, March 2026 \\
Local LLM 1.5B (Qomni) & 8,000 & 45,000 & qwen2.5-1.5b-instruct-q4 \\
\bottomrule
\end{tabular}
\end{table}

\subsection{Measured CRYS-L Benchmark (Server5, 2026-04-13)}

\begin{table}[h]
\centering
\caption{5-run CRYS-L benchmark — Server5 (AMD EPYC 12-core)}
\begin{tabular}{lrrrrl}
\toprule
Plan & Run1 & Run2 & Run3 & Run4 & Run5 & Avg \\
\midrule
pump\_sizing (500gpm,80psi) & 302 & 277 & 305 & 282 & 265 & \textbf{286ms} \\
hazen\_williams (150mm,300m) & 56 & 51 & 61 & 50 & 49 & \textbf{53ms} \\
voltage\_drop (75kw,380v,120m) & 130 & 118 & 113 & 127 & 112 & \textbf{120ms} \\
\midrule
\multicolumn{6}{l}{Overall avg} & \textbf{153ms} \\
\bottomrule
\end{tabular}
\end{table}

\textbf{Key observation}: The 286ms pump\_sizing result uses the Qomni Planner
(crystal retrieval + oracle), while hazen\_williams at 53ms uses the
Symbolic Engine (pure formula, no retrieval). The WASM execution itself
takes $< 1$~ms in both cases; the overhead is HTTP + JSON + retrieval.

\subsection{Accuracy Comparison}

\begin{table}[h]
\centering
\caption{Accuracy comparison — pump sizing (500 GPM, 80 PSI, $\eta=0.75$)}
\begin{tabular}{lllll}
\toprule
Approach & Answer & Exact? & Standard cited? & Reproducible? \\
\midrule
\textbf{CRYS-L JIT} & \textbf{13.26 HP} & \textbf{Yes} & \textbf{NFPA 20:2022 §4.26} & \textbf{Yes} \\
GPT-4 Turbo & ``about 15 HP'' & No ($\pm$15\%) & No & No \\
Claude 3.5 Sonnet & ``13-14 HP'' & Partial & No & No \\
Python (manual) & 13.26 HP & Yes & Manual & Yes \\
Excel (manual) & 13.26 HP & Yes & Manual & Yes \\
Local LLM 1.5B & ``12-18 HP'' & No ($\pm$25\%) & No & No \\
\bottomrule
\end{tabular}
\end{table}

The ground truth is: $\text{HP} = (500 \times 0.06309 \times 80 \times 0.70307) / (0.75 \times 76.04) = 13.26$~HP.

LLMs introduce 15--25\% error on this calculation because they pattern-match
the formula from training data rather than executing it algebraically.
CRYS-L, Python, and Excel all produce the exact answer, but only CRYS-L
does so without manual setup and with standard citation.

\subsection{Qualitative Comparison}

\begin{table}[h]
\centering
\caption{Qualitative comparison across tools}
\begin{tabular}{lcccccc}
\toprule
& CRYS-L & GPT-4 & Claude & Python & Excel & Eng.~SW \\
\midrule
Exact answer & \checkmark & $\times$ & $\times$ & \checkmark & \checkmark & \checkmark \\
Standard cited & \checkmark & $\times$ & $\times$ & Manual & Manual & \checkmark \\
No setup needed & \checkmark & \checkmark & \checkmark & $\times$ & $\times$ & $\times$ \\
Works offline & \checkmark & $\times$ & $\times$ & \checkmark & \checkmark & \checkmark \\
Embeddable (API) & \checkmark & \checkmark & \checkmark & Partial & $\times$ & $\times$ \\
Open / extensible & \checkmark & $\times$ & $\times$ & \checkmark & $\times$ & $\times$ \\
Cost & Free & \$0.01/call & \$0.015/call & Free & License & \$500+ \\
Latency & 53-286ms & 3-8s & 2-7s & 80-300ms & 50-200ms & Instant \\
\bottomrule
\end{tabular}
\end{table}

\subsection{When LLMs Win}

CRYS-L is not a replacement for LLMs on unstructured or creative tasks.
LLMs outperform CRYS-L when:
\begin{itemize}
  \item The query requires interpretation of ambiguous requirements
  \item The domain has no formal standard (soft design decisions)
  \item The user needs explanation, not just a number
  \item The calculation is novel and not in the CRYS-L stdlib
\end{itemize}

The optimal architecture (implemented in Qomni Engine) uses both:
CRYS-L handles the deterministic calculation layer, the LLM handles
explanation and edge cases.

%─────────────────────────────────────────────────────────────────────────────
\section{Community: How to Contribute}

\subsection{Writing a New Plan}

A valid CRYS-L plan requires:
\begin{enumerate}
  \item A \texttt{meta\{\}} block with \texttt{standard}, \texttt{source}, \texttt{domain}
  \item At least one \texttt{assert} per input parameter
  \item At least one \texttt{output} with label and unit
  \item A \texttt{formula} declaration for each key equation
\end{enumerate}

Supported domains for new contributions:

\begin{table}[h]
\centering
\caption{CRYS-L v2 domain registry}
\begin{tabular}{lll}
\toprule
Domain & Reference Standards & Example Plans \\
\midrule
\texttt{hidraulica} & IS.010/IS.020 Peru, AWWA M22 & hazen\_williams, darcy\_weisbach \\
\texttt{nfpa\_electrico} & NFPA 13/20/72, IEC 60364 & pump\_sizing, voltage\_drop \\
\texttt{civil} & ACI 318, NTE Peru & beam\_deflection, column\_load \\
\texttt{mecanica} & ISO 6336, AGMA & shaft\_power, gear\_ratio \\
\texttt{termica} & ASHRAE 90.1, EN ISO 6946 & hvac\_load, insulation \\
\texttt{sanitaria} & IS.010 Peru, RNE OS.050 & water\_demand, tank\_sizing \\
\bottomrule
\end{tabular}
\end{table}

\subsection{Test Vector Format}

Each plan PR must include a test file:

\begin{lstlisting}[caption=Test vector format — datasets/tests/plan\_name.json]
{
  "plan": "plan_hazen_williams",
  "cases": [
    {
      "inputs": {"Q": 2.5, "D": 150.0, "C": 120.0, "L": 200.0},
      "expected": {"V": 0.141, "h_f": 0.089},
      "tolerance_pct": 0.5,
      "source": "IS.010 Peru, Example 3.2, Table A"
    }
  ]
}
\end{lstlisting}

\subsection{Contribution Checklist}

\begin{itemize}
  \item[$\square$] \texttt{meta\{\}} block with standard + section reference
  \item[$\square$] All inputs have \texttt{assert} with meaningful message
  \item[$\square$] All key formulas have \texttt{formula} declaration string
  \item[$\square$] All outputs have \texttt{label} and \texttt{unit}
  \item[$\square$] Test vector file with at least 3 cases from published tables
  \item[$\square$] Plan name follows \texttt{plan\_\{domain\}\_\{calc\}} convention
\end{itemize}

%─────────────────────────────────────────────────────────────────────────────
\section{Live Demo}

CRYS-L v2.1 is live and accessible at:
\begin{center}
  \href{https://qomni.clanmarketer.com/crysl/}{\textbf{qomni.clanmarketer.com/crysl/}}
\end{center}

The demo runs WASM entirely in the browser — no server round-trip for execution.
Try:
\begin{itemize}
  \item \texttt{plan\_pump\_sizing(500, 80, 0.75)} — NFPA 20 fire pump
  \item \texttt{plan\_hazen\_williams(3.0, 150, 120, 200)} — Hazen-Williams
  \item \texttt{plan\_voltage\_drop(75, 380, 120)} — IEC cable sizing
\end{itemize}

%─────────────────────────────────────────────────────────────────────────────
\section{Conclusion}

CRYS-L v2 provides a principled middle ground between the rigidity of
commercial engineering software and the imprecision of LLM-based calculation.
It achieves:
\begin{itemize}
  \item \textbf{53--286~ms} end-to-end latency (15--850$\times$ faster than LLM)
  \item \textbf{Exact answers} from published engineering standards
  \item \textbf{Zero cost} inference (no API calls, no GPU)
  \item \textbf{Portable WASM} runtime for browser, server, and embedded
  \item \textbf{Open standard} with community PR path
\end{itemize}

The CRYS-L language, standard library, runtime, and this paper are released
under MIT license at \href{https://github.com/condesi/crysl-lang}{condesi/crysl-lang}.

\bibliographystyle{plain}
\begin{thebibliography}{9}
\bibitem{nfpa20}
  NFPA, ``NFPA 20: Standard for the Installation of Stationary Pumps
  for Fire Protection,'' National Fire Protection Association, 2022.
\bibitem{iec60364}
  IEC, ``IEC 60364-5-52: Electrical Installations — Selection and Erection
  of Electrical Equipment,'' International Electrotechnical Commission, 2009.
\bibitem{is010}
  Ministerio de Vivienda Peru, ``IS.010: Instalaciones Sanitarias para
  Edificaciones,'' Reglamento Nacional de Edificaciones, 2006.
\bibitem{wasmtime}
  Bytecode Alliance, ``Wasmtime: A Fast and Secure Runtime for WebAssembly,''
  \href{https://wasmtime.dev}{wasmtime.dev}, 2024.
\bibitem{qomni2026}
  P.~Rojas Masgo and Qomni AI, ``Qomni Engine v8: From Crystal Distillation
  to Adaptive Cognitive OS,'' \emph{Preprint}, April 2026.
\end{thebibliography}

\end{document}
